\documentclass[10pt,a4paper]{amsart}
\usepackage[margin=19mm]{geometry}
\usepackage{amsmath,amssymb,mathtools}
\usepackage[scr=rsfs,cal=euler]{mathalfa}
\usepackage{xfrac}
\usepackage[unicode,hidelinks,bookmarks=false]{hyperref}
\newcommand{\ie}{i.e.\ }
\newcommand{\cf}{cf.\ }
\newcommand{\resp}{respectively\ }
\newcommand{\Subsection}{\subsection}
\providecommand{\nobreakdash}{\nobreak\hbox{-}}
\setlength{\emergencystretch}{2em}
\multlinegap0pt
\usepackage{amssymb}
\usepackage{xcolor}
\usepackage{mathtools}
\usepackage{float}
\newtheorem{theorem}{Theorem}
\title{Polynomial and spectra factorization of graphs obtained by iteration the operad of generalized graph composition.}
\author{Jean Liendo}
\date{Source: arXiv:2604.27097v1 (CC BY 4.0)}
\begin{document}
\maketitle

\noindent\emph{Source and licence.} Polynomial and spectra factorization of graphs obtained by iteration the operad of generalized graph composition.. Version arXiv:2604.27097v1 (CC BY 4.0), distributed by arXiv under the Creative Commons Attribution 4.0 International licence, http://creativecommons.org/licenses/by/4.0/, as recorded in the arXiv licence metadata for this exact version and shown on the abstract page https://arxiv.org/abs/2604.27097v1. Full licence notice: \texttt{licenses/arxiv-2604.27097-CC-BY-4.0.txt}.\par\smallskip

\noindent\textit{Editorial scope.} Counted unit: the article's theorem iterated (source0.tex lines 485--490). The article's complete original proof of it is delivered with it (source0.tex lines 492--511, one \texttt{proof} environment of 20 lines). No mathematics is written, repaired or re-derived: the statement and the proof are the article's own lines, in the article's own notation, and no retained line is substituted or re-flowed.  Retained uncounted as the source-local context the proof needs: lines 479--481.  Nothing was omitted from any retained span, so the package carries no omission record.  No retained line contains a citation, so the wrapper carries no bibliography and no bibliography entry.  No retained line contains a figure environment, an image-inclusion command or drawing code, so no image is delivered and the package declares no asset dependency.  Editorial formatting only, in addition: an AMS \texttt{amsart} preamble that restates the article's own declarations and macros used by the retained lines, namely the environments \texttt{theorem} on separate counters, together with the standard packages those lines need.  The retained environments print in this document as Theorem 1 in order of appearance, whereas the article prints the same items with the numbering of its own full document.  The complete native source of the article as distributed in the arXiv e-print for this exact version is bundled unchanged as \texttt{sources/199444/source0.tex}.
\begin{equation}\label{espectro adyacencia Join arbol}
\sigma\left(A\left(\bigvee_{g_v} a_v\right)\right)=\dfrac{\sigma(A(a_v))}{\mathrm{reg}(a_v)}\cdot \sigma(A(a_v, g_v)).
\end{equation}

\begin{theorem}\label{adjacency spectrum iterated}
Let $\mathscr{T}\in \mathscr{F}_{\mathscr{G}_+}[\ell]$ be a factorization of a graph $g$, that is, $\displaystyle{\bigvee \mathscr{T}=g}$. Suppose that for each internal vertex $v$ of $\mathscr{T}$, $a_v$ is an ordered assembly of regular graphs. Then
\begin{equation}\label{espectro adyacencia sobre un arbol enriquecido}
\sigma(A(g))=0^{|\ell|}\cdot \prod_{v\in \mathrm{Iv}(\mathscr{T})}\dfrac{\sigma(A(a_v,g_v))}{\mathrm{reg}(a_v)}
\end{equation}
\end{theorem}

\begin{proof}
Let $\psi$ be the function given by
\[A\left(\bigvee \mathscr{T}\right) \mapsto \displaystyle{0^{|\ell|}\cdot \prod_{v\in \mathrm{Iv}(\mathscr{T})}\dfrac{\sigma(A(a_v,g_v))}{\mathrm{reg}(a_v)}}.\]

Let $r$ be the root of $\mathscr{T}$, let $s(r)=\{r_1, r_2,...,r_k\}$ be the set of childrens of $r$. Then, for each $i=1,2,3,...,k$, we have
\begin{equation*}
\dfrac{\psi\left(A\left(\bigvee \mathscr{T}_{r_i}\right)\right)}{\mathrm{reg}\left(\bigvee \mathscr{T}_{r_i}\right)}=\left(0^{|\ell_{r_i}|}\cdot \prod_{w\in \mathrm{Iv}(\mathscr{T}_{r_i})}\dfrac{\sigma(A(a_w,g_w))}{\mathrm{reg}(a_w)}\right)\cdot \dfrac{1}{\mathrm{reg}\left(\bigvee \mathscr{T}_{r_i}\right)}.
\end{equation*}
Therefore,
\begin{eqnarray*}
\prod^k_{i=1} \dfrac{\psi\left(A\left(\bigvee \mathscr{T}_{r_i}\right)\right)}{\mathrm{reg}\left(\bigvee \mathscr{T}_{r_i}\right)}&=&0^{|\ell|}\cdot \left(\prod_{w\in \mathrm{Iv}(\mathscr{T}) \atop w\neq r} \dfrac{\sigma(A(a_w,g_w))}{\mathrm{reg}(a_w)}\right)\cdot \dfrac{1}{\mathrm{reg}(a_r)}
\end{eqnarray*}

which implies that
\begin{eqnarray*}
\dfrac{\psi(A(a_r))}{\mathrm{reg}(a_r)}\cdot \sigma(A(a_r, g_r))&=&0^{|\ell|}\left(\prod_{w\in \mathrm{Iv}(\mathscr{T})\atop w\neq r} \dfrac{\sigma(A(a_w,g_w))}{\mathrm{reg}(a_w)}\right)\cdot \dfrac{\sigma(A(a_r,g_r))}{\mathrm{reg}(a_r)}\\
&=&\psi\left(A\left(\bigvee \mathscr{T}\right)\right).
\end{eqnarray*}
In conclusion, $\psi$ satisfies the recursive equation (\ref{espectro adyacencia Join arbol}), so $\psi=\sigma$.
\end{proof}

\end{document}
